\begin{textsample}{POS Dim 2 – global\_south\_2025\_02 – Score 12.00 – global\_south\_2025\_02/121796245.txt}  \label{ex:f2_pos_374}
Israeli army soldiers stand \textbf{guard} as Israeli Jewish settlers tour the old market in the city of Hebron in the occupied West Bank AFP The probe report by the Israeli Defence Forces ( IDF ) on October 7 revealed how it misread intelligence material for years, believing the \textbf{militant} group Hamas in Gaza did not pose a significant threat to Israel.
It also believed Hamas leader Yahya Sinwar, who Israel killed in a \textbf{drone} \textbf{attack}, was a pragmatist who didn't want a war with Israel.
The top-level investigation into the military ' s failures during the lead-up to the \textbf{attack} was released on Thursday.
Here are some important findings presented by the IDF : 1 ) The IDF ' s Gaza division was"defeated"for several hours as chaos and confusion slowed the fightback.
There were just 767 IDF troops at the \textbf{border} for over 5,000 terrorists on that day. 2 ) The General Staff did not realise that the Gaza Division had fallen or the severity of the \textbf{attack}, thereby failing to put @ @ @ @ @ @ @ @ @ @ 3 ) Most of Hamas ' s atrocities were carried out within the first six hours of the \textbf{attack} and during those hours, many Israeli \textbf{border} communities, IDF posts, and main routes in the western Negev area were under Hamas control. 4 ) The intelligence material was insistently misinterpreted over the years and the military overly relied on an early warning to prepare its defences. 5 ) The IDF assumed there was no imminent threat and that the military ' s deployment on the \textbf{border} was according to protocol. 6 ) The IDF believed the Hamas terror group in Gaza did not pose a significant threat to Israel and that it was uninterested in a large-scale war.
The military assumed that Hamas saw its civil control in Gaza as a strategic asset and wanted to \textbf{strike} a deal with Israel. 6 ) Israel believed that its tunnel networks had been degraded and that any \textbf{cross-border} threat would be thwarted by Israel ' s high-tech \textbf{border} fence. 7 ) Though the Military Intelligence Directorate had intel outlining Hamas ' s intent @ @ @ @ @ @ @ @ @ @ dismissed as unrealistic and unfeasible. 8 ) The Military Intelligence Directorate believed that Hamas leader Yahya Sinwar was a pragmatist who was not seeking a major \textbf{escalation} with Israel. ( 9 ) Israeli military officials believed they had intelligence superiority over Hamas and assumed that Hamas could never surprise them due to the abundant information they had about the \textbf{militant} group. 10 ) The IDF believed Hamas was also carrying out minimal efforts to prepare for \textbf{infiltrations} and \textbf{drone} \textbf{attacks} and was focusing on building up its \textbf{rocket} \textbf{launching} \textbf{capabilities}.

% matched lemmas: attack, border, capability, cross-border, drone, escalation, guard, infiltration, launch, militant, rocket, strike
\end{textsample}
